%% Late Binding in Machine-Consumed Instruction Files
%% ACM TOSEM submission skeleton — values from frozen exports only.
%% Compile: pdflatex main && bibtex main && pdflatex main && pdflatex main

\documentclass[acmsmall,review,anonymous]{acmart}

\usepackage{xcolor}
\newcommand{\todo}[1]{\textcolor{red}{[TODO: #1]}}

\AtBeginDocument{%
  \providecommand\BibTeX{{Bib\TeX}}}

\setcopyright{none}
\copyrightyear{2026}
\acmYear{2026}
\acmDOI{TODO}
\acmJournal{TOSEM}
\acmVolume{TODO}
\acmNumber{TODO}
\acmArticle{TODO}

\title{Late Binding in Machine-Consumed Instruction Files}

\author{TODO Author One}
\affiliation{%
  \institution{TODO Institution}
  \city{TODO}
  \country{TODO}}
\email{todo@example.com}

\author{TODO Author Two}
\affiliation{%
  \institution{TODO Institution}
  \city{TODO}
  \country{TODO}}
\email{todo@example.com}

\begin{abstract}
Machine-consumed instruction files (\eg AGENTS.md, skills manifests, IDE rule packs) are increasingly treated as project documentation, yet coding agents consume them under different semantics than human readers.
We propose a \emph{late-binding model}: these artifacts combine a \textbf{directive channel} that shapes agent behavior with a \textbf{referential channel} whose pointers are resolved against the live repository at use time.
Using a longitudinal panel of instruction-file references (RQ1--RQ4), confirmatory audits of born-stale claims, validated agent-maintenance attribution (P4), and a controlled A/B/C agent experiment (RQ5), we show that (i) genuine post-verification reference decay is rare in our cohort (0/121 adjusted failures), (ii) a large class of references is \emph{confirmed false at creation} (1200/1405; 85.4\%), (iii) agents read instructions universally when present (128/128) and follow anchor references in a majority of runs (77.8\% on condition~B), (iv) instruction presence amplifies edit scope ($\sim$100 vs $\sim$2 files modified), yet (v) static referential truth is weakly coupled to task success (paired $\Delta$ A$-$B = 0.00~pp on 63 triplets).
We argue that observational \emph{truth decay} and causal \emph{operational cost} must be separated, and we document threats to validity and unsupported claims explicitly.
\end{abstract}

\begin{CCSXML}
<ccs2012>
<concept>
<concept_id>10011007.10011006</concept_id>
<concept_desc>Software and its engineering~Software maintenance tools</concept_desc>
<concept_significance>500</concept_significance>
</concept>
<concept>
<concept_id>10010147.10010178</concept_id>
<concept_desc>Computing methodologies~Artificial intelligence</concept_desc>
<concept_significance>300</concept_significance>
</concept>
</ccs2012>
\end{CCSXML}

\ccsdesc[500]{Software and its engineering~Software maintenance tools}
\ccsdesc[300]{Computing methodologies~Artificial intelligence}

\keywords{software documentation, coding agents, instruction files, reference integrity, empirical study}

\begin{document}
\maketitle

\section{Introduction}
\label{sec:introduction}

Coding agents consume repository-local instruction files---\texttt{AGENTS.md}, Cursor rules, Claude skills, Copilot instruction packs---as authoritative context at task time.
Practitioners and researchers often analyze these files as if they were human-facing documentation whose truth is fixed at commit time.
We argue that this framing is insufficient: machine-consumed instruction files are \textbf{late-binding artifacts} whose semantics split across two channels.

\paragraph{Central claim.}
Machine-consumed instruction files are not ordinary documentation.
They operate through (1) a \emph{directive channel} that strongly shapes agent behavior and (2) a \emph{referential channel} whose \emph{static truth} is only weakly coupled to task outcomes because references are resolved against the live repository at use time.

\paragraph{Empirical program.}
We report a multi-stage empirical program on a frozen laboratory corpus:
longitudinal reconstruction and survival analysis (RQ1--RQ2),
observational integrity and lifecycle dynamics (RQ3--RQ4),
confirmatory audits of born-stale false claims,
human-validated agent-maintenance attribution (P4),
and a controlled agent experiment comparing truthful, confirmed-false, and no-instruction conditions (RQ5 A/B/C).

\paragraph{Contributions.}
\begin{itemize}
  \item A conceptual \emph{late-binding model} with explicit constructs (directive vs.\ referential channels, static truth, runtime resolution, causal load-bearing references).
  \item Observational evidence that post-verification reference decay is negligible in our audited cohort (0/121 adjusted genuine decay) while born-false claims at creation are prevalent (1200/1405 confirmed false; 85.4\%).
  \item Agent trace evidence that instructions are read (128/128) and often followed (72.3--77.8\% by condition), amplifying edit scope ($\sim$100 vs $\sim$2 files modified), without a measurable success difference between truthful and false referential content on paired triplets ($\Delta$ = 0.00~pp; McNemar $p = 1.0$).
  \item A reproducibility bundle mapping each claim to frozen exports (\path{exports/paper_synthesis/late_binding_evidence_table.csv}).
\end{itemize}

\paragraph{Paper organization.}
\todo{One-sentence roadmap.}

\section{Background}
\label{sec:background}

\subsection{Machine-consumed instruction files}
\todo{Define AGENTS.md, skills, IDE rules; contrast with README/CONTRIBUTING (P5 descriptive baseline). Cite practitioner docs.}

\subsection{Reference integrity and truth decay}
Longitudinal panels label verifiable references as \textsc{Verified} or \textsc{Missing} at each commit.
Prior work in this laboratory distinguishes:
\begin{itemize}
  \item \textbf{Born-stale} references: \textsc{Missing} without a prior \textsc{Verified} observation.
  \item \textbf{Post-verification decay}: transition from \textsc{Verified} to \textsc{Missing}.
\end{itemize}
RQ3 explicitly warns against merging these prevalence rates.

\subsection{Agent maintenance attribution}
Gate P4 validates whether commits reflect genuine agent-assisted instruction maintenance.
Human gold validation on $N=200$ commits yields precision 0.958 and F1 0.9617 for the positive class \texttt{true\_agent\_maintenance} (\path{exports/truth_pilot/p4_validation.md}).

\subsection{Controlled agent impact experiments}
RQ5 manipulates instruction-file content at a pinned commit while holding repository state, task prompt, and test command fixed.
The redesigned protocol adds condition~C (no instruction file) to separate instruction \emph{presence} from referential \emph{truth}.

\section{Conceptual Model}
\label{sec:model}

We adopt the late-binding model documented in \path{docs/LATE_BINDING_MODEL_v1.md}.
Figure~\ref{fig:model-dag} summarizes the causal structure; Table~\ref{tab:constructs} defines core constructs.

\begin{figure}[t]
  \centering
  % Source: docs/LATE_BINDING_MODEL_v1.md (Mermaid DAG); export to PDF for camera-ready.
  \fbox{\parbox{0.92\linewidth}{\centering\vspace{2em}
    \textbf{TODO:} Include model DAG figure\\
    (\path{paper/figures/model-dag.pdf})\\
    \vspace{2em}}}
  \caption{Late-binding model: static panel truth, two-channel consumption, environmental task difficulty.}
  \label{fig:model-dag}
\end{figure}

\begin{table}[t]
  \caption{Core constructs in the late-binding model.}
  \label{tab:constructs}
  \begin{tabular}{@{}p{0.22\linewidth}p{0.70\linewidth}@{}}
\toprule
\textbf{Construct} & \textbf{Definition} \\
\midrule
Directive channel & Normative text shaping agent behavior without resolvable pointers \\
Referential channel & Pointers resolved against live repository at use time \\
Late binding & Resolution deferred to agent consumption, not authoring \\
Static truth & Mechanical VERIFIED/MISSING at pinned commit \\
Runtime resolution & Agent read/follow/act on referential pointers \\
Read--repair & Agent detects stale pointer and substitutes without external help \\
Causal load-bearing reference & Pointer on task-critical path affecting outcome \\
Environmental task difficulty & Repo-specific test/build friction gating success \\
\bottomrule
\end{tabular}

\end{table}

\paragraph{Directive channel.}
Normative text shapes agent goals and tool policy without requiring a resolvable pointer (\eg ``run tests before finishing'').
RQ5 shows this channel activates even when the referential channel is absent: mean files modified drops from $\sim$100 (A/B) to 1.7 (C) on 63 paired triplets.

\paragraph{Referential channel.}
Pointers (\paths, commands, configs) intended for runtime resolution.
Static truth is assessed mechanically in RQ1--RQ4; runtime resolution is observed in RQ5 traces (\texttt{instruction\_read}, \texttt{instruction\_followed}).

\paragraph{Late binding.}
Binding occurs at agent consumption time against the workspace snapshot, not at authoring time.
A statically false pointer may still be \emph{read} and \emph{followed} (77.8\% of condition~B runs in uptake analysis).

\paragraph{Causal load-bearing vs.\ peripheral references.}
Selection effects (85.4\% of matched pairs: cited path churn $\leq$ uncited control) suggest many cited paths are intrinsically stable; false claims may be non-load-bearing with respect to task success.

\section{Study Design}
\label{sec:design}

\subsection{Corpus and gates}
\todo{Summarize registry, P1--P5 gates, RQ1 longitudinal export path.}
All reported numbers derive from frozen exports under \path{exports/truth_decay_pilot/}, \path{exports/truth_pilot/}, \path{exports/rq5_agent_impact/}, and \path{exports/rq5_agent_impact_c/}.

\subsection{RQ1--RQ4: Observational track}
\begin{itemize}
  \item \textbf{RQ1:} Longitudinal feasibility --- reference states over commit history.
  \item \textbf{RQ2:} Survival / first-missing events after verification.
  \item \textbf{RQ3:} Observational integrity by maintenance regime.
  \item \textbf{RQ4:} Multi-state lifecycle dynamics.
\end{itemize}
\todo{Insert cohort sizes from \path{rq1_feasibility.md}, \path{rq2_summary.md}.}

\subsection{Confirmatory audits}
\textbf{Born-stale autopsy} classifies 17{,}747 born-stale references into heterogeneous categories (\path{born_stale_summary.md}).
\textbf{GFC confirmatory audit} re-evaluates 1{,}405 \texttt{genuine\_false\_claim} labels.
\textbf{RQ2 failure audit} adjudicates 121 post-verification \texttt{first\_missing} events.
\textbf{Cited vs.\ uncited audit} matches 2{,}259 path pairs across 78 repositories.

\subsection{P4: Agent maintenance validation}
Human gold labels on 200 commits; frozen heuristic classifier evaluated without retraining.

\subsection{RQ5: Controlled agent experiment}
\begin{itemize}
  \item \textbf{A:} Truthful instruction blob at pinned commit.
  \item \textbf{B:} Confirmed-false instruction blob; all else identical.
  \item \textbf{C:} Instruction file removed (no-instruction baseline).
\end{itemize}
Agent: Claude Code CLI (\texttt{claude\_code} adapter).
Partial A/B: 128 runs on 22/35 cases (\path{exports/rq5_agent_impact/rq5_results.csv}).
Complete C: 105 runs on 35/35 cases (\path{exports/rq5_agent_impact_c/rq5_results.csv}).
Paired ABC analysis: 63 triplets (22-case overlap).
\todo{Insert exact task prompt and test-command policy from protocol.}

\section{Results}
\label{sec:results}

We organize results by research question.
All values below are copied from frozen exports; see Table~\ref{tab:evidence-map} and \path{exports/paper_synthesis/late_binding_evidence_table.csv}.

\begin{table*}[t]
  \caption{Selected frozen evidence mapped to paper claims (full table in CSV).}
  \label{tab:evidence-map}
  \small
  \begin{tabular}{@{}p{0.28\linewidth}p{0.22\linewidth}p{0.14\linewidth}p{0.26\linewidth}@{}}
  \toprule
  \textbf{Claim} & \textbf{Source export} & \textbf{Value} & \textbf{Limitation} \\
  \midrule
  Zero genuine post-verification decay & \path{rq2_failure_audit_summary.md} & 0/121 & Single audit pass \\
  Confirmed false at creation & \path{gfc_confirmatory_summary.md} & 1200/1405 (85.4\%) & Wilson CI 83.5--87.2\% \\
  Cited paths stable vs.\ controls & \path{cited_uncited_summary.md} & 85.4\% pairs & Mean churn CI crosses 0 \\
  Agent maintenance F1 & \path{p4_validation.md} & 0.9617 & $N=200$ gold sample \\
  Instruction read (A/B) & \path{rq5_uptake_analysis.md} & 128/128 & Claude Code only \\
  Instruction followed (B) & \path{rq5_uptake_analysis.md} & 77.8\% (49/63) & Trace heuristic \\
  Paired $\Delta$ success A$-$B & \path{rq5_abc_comparative_analysis.md} & 0.00 pp & 63 triplets \\
  Files modified A vs.\ C & \path{rq5_abc_comparative_analysis.md} & 99.7 vs.\ 1.7 & Behavioral, not outcome \\
  Condition C complete & \path{rq5_agent_impact_c/} & 105/105 runs & No matched A/B on 13 cases \\
  \bottomrule
  \end{tabular}
\end{table*}


\subsection{RQ1--RQ2: Static reference integrity}
\paragraph{Post-verification decay.}
Of 121 post-verification failure events audited, adjusted genuine decay count is \textbf{0/121} (0.0\%; Wilson 95\% CI: 0.0\%--3.1\%) (\path{rq2_failure_audit_summary.md}).
73.6\% of failures are classified as extractor artifacts.

\paragraph{Born-stale prevalence.}
Raw \texttt{genuine\_false\_claim} count: 1{,}405 (7.9\% of 17{,}747 born-stale cohort) (\path{born_stale_summary.md}).

\subsection{Confirmatory audits}
\paragraph{Confirmed false at creation.}
GFC confirmatory audit: \textbf{1{,}200 / 1{,}405} (85.4\%; Wilson 95\% CI: 83.5\%--87.2\%) (\path{gfc_confirmatory_summary.md}).
Adjusted born-stale false-claim rate: 6.76\% of cohort.

\paragraph{Path stability selection.}
Cited vs.\ uncited audit: in \textbf{85.4\%} of 2{,}259 matched pairs, cited path git churn $\leq$ uncited control (\path{cited_uncited_summary.md}).
Mean paired difference in churn: 0.44 commits (95\% bootstrap CI: $-0.21$--1.15; includes zero).

\subsection{P4: Agent maintenance attribution}
Precision 0.958; F1 0.9617 on $N=200$ human-reviewed commits (\path{p4_validation.md}).

\subsection{RQ5: Agent behavior and outcomes}

\subsubsection{Uptake funnel (A/B; $n=128$)}
\begin{itemize}
  \item \texttt{instruction\_read}: 128/128 (100\%).
  \item \texttt{instruction\_followed}: 47/65 (72.3\%) condition A; 49/63 (77.8\%) condition B.
  \item Task success: 8/65 (12.3\%) A; 8/63 (12.7\%) B.
\end{itemize}
Source: \path{rq5_uptake_analysis.md}.

\subsubsection{Behavioral amplification (paired triplets; $n=63$)}
Mean files modified: A = 99.7; B = 103.0; C = 1.7 (\path{rq5_abc_comparative_analysis.md}).
Compilation success: 100\% all conditions on triplets.

\subsubsection{Causal contrasts (paired triplets)}
\begin{itemize}
  \item \textbf{A $-$ B:} $\Delta$ success = 0.00~pp; cluster bootstrap 95\% CI [$-6.35$, $+6.35$]; McNemar $p = 1.0$.
  \item \textbf{A $-$ C:} $\Delta$ success = +4.76~pp; cluster bootstrap 95\% CI [$-12.70$, $+22.22$]; McNemar $p = 0.5811$.
  \item \textbf{B $-$ C:} $\Delta$ success = +4.76~pp (identical to A$-$C on overlap triplets where A and B success patterns match).
\end{itemize}
Success rates on triplets: A = 12.7\% (8/63); B = 12.7\% (8/63); C = 7.9\% (5/63).

\subsubsection{Condition C complete cohort ($n=105$)}
Aggregate success: 12/105 (11.4\%) (\path{exports/rq5_agent_impact_c/rq5_results.csv}).

\paragraph{Figures.}
\todo{Reference paper/figures copies of export PDFs for camera-ready.}

\section{Discussion}
\label{sec:discussion}

\subsection{Reconciling observational and causal evidence}
Observational work establishes widespread \emph{static} referential problems (born-false prevalence, heterogeneous taxonomy) while showing negligible \emph{genuine post-verification decay} in the audited cohort.
RQ5 shows agents \emph{consume} instruction files (100\% read) and often \emph{bind} to anchor references ($\sim$77\% on B) yet \emph{do not} exhibit different success rates when static referential truth is swapped (A vs.\ B).

The late-binding model reconciles these findings:
static panel labels measure authoring-time alignment; task outcomes are gated primarily by \textbf{environmental task difficulty} ($\sim$88\% failure on unsuccessful runs) and secondarily by whether a reference is \textbf{causally load-bearing}.

\subsection{Directive vs.\ referential coupling}
The $\sim$50$\times$ gap in files modified between instruction-present (A/B) and absent (C) conditions indicates the \textbf{directive channel is strongly executive}: it expands agent edit scope even when success remains low.
Referential static truth, by contrast, appears \textbf{weakly coupled} to success in our pilot tasks.

\subsection{Implications for ``Truth Debt''}
Operational cost claims require measurable downstream harm.
Current RQ5 evidence \textbf{does not support} promoting Truth Debt beyond observational Truth Decay: paired A$-$B success difference is null on available triplets.
\todo{Qualify Truth Debt paragraph with CI width and power limits.}

\subsection{Alternative explanations considered}
\begin{itemize}
  \item \textbf{Decorative instructions:} weakened by 100\% read and majority follow rates.
  \item \textbf{Agent robustness to false refs:} weakened; false claim used in 77.8\% of B runs without improved success.
  \item \textbf{Measurement artifact only:} partially supported for post-verification decay; not for confirmed born-false claims.
  \item \textbf{Task difficulty dominates:} supported by high \texttt{tests\_failed} counts and 100\% compile success.
\end{itemize}

\subsection{Implications for tool builders}
Separate linting for directive vs.\ referential content; resolve pointers at agent start; annotate load-bearing references; telemetry for read--repair events.
See \path{docs/LATE_BINDING_MODEL_v1.md} \S7.

\section{Threats to Validity}
\label{sec:threats}

\paragraph{Construct validity.}
\texttt{instruction\_followed} is inferred from tool traces, not ground-truth binding.
VERIFIED/MISSING states depend on extraction heuristics (73.6\% extractor-artifact rate among post-verification failures).

\paragraph{Internal validity.}
RQ5 ABC overlap covers 22/35 cases (63 triplets); 13 cases have C only.
A/B and C runs live in separate export directories; pairing assumes matching \texttt{case\_id} and \texttt{replicate\_id}.

\paragraph{External validity.}
Single agent platform (Claude Code CLI).
Corpus drawn from open-source AI-convention repositories; \todo{state registry size}.

\paragraph{Statistical conclusion validity.}
Low success rates ($\sim$11--13\%) yield wide confidence intervals; cluster bootstrap CIs for A$-$C and B$-$C include zero.
McNemar tests are underpowered at $n=63$ triplets.

\paragraph{Ecological validity.}
Pinned-commit workspaces with injected instruction blobs may differ from IDE-integrated agent sessions.

\paragraph{Supported vs.\ unsupported claims.}
Table~\ref{tab:evidence-map} maps evidentiary limits per claim.
Section~8 of \path{docs/LATE_BINDING_MODEL_v1.md} lists claims \emph{not} supported.

\section{Data and Reproducibility}
\label{sec:reproducibility}

\subsection{Artifact structure}
\begin{itemize}
  \item \path{exports/truth_decay_pilot/} --- RQ1--RQ4 tables, audits, figures.
  \item \path{exports/truth_pilot/p4_validation.md} --- attribution validation.
  \item \path{exports/rq5_agent_impact/} --- partial A/B runs, uptake, ABC analysis.
  \item \path{exports/rq5_agent_impact_c/} --- complete condition~C runs.
  \item \path{exports/paper_synthesis/late_binding_evidence_table.csv} --- claim-to-evidence map.
  \item \path{docs/LATE_BINDING_MODEL_v1.md} --- conceptual synthesis.
  \item \path{protocol/RQ5_AGENT_IMPACT_EXPERIMENT_v1.md} --- experiment protocol v1.1.
\end{itemize}

\subsection{Replication commands}
\todo{Insert exact Makefile targets from repository freeze commit.}
Example frozen invocations (no agent reruns required for observational track):
\begin{verbatim}
make truth-decay-rq1 truth-decay-rq2 ...
make truth-decay-rq5-redesign   # plan only
\end{verbatim}

\subsection{Experiment freeze statement}
As of paper-writing mode, \textbf{no further agent experiments are scheduled}.
RQ5 A/B partial (128 runs) and C complete (105 runs) are frozen.

\subsection{Evidence table maintenance}
Each row in \path{late_binding_evidence_table.csv} includes a \texttt{limitation} column.
Authors must not cite a row without its stated limitation in the narrative.

\section{Conclusion}
\label{sec:conclusion}

Machine-consumed instruction files behave as late-binding artifacts: a directive channel shapes agent behavior strongly, while a referential channel carries static claims whose truth is evaluated separately from runtime resolution.
Across a frozen multi-RQ laboratory program, we find abundant born-false referential content at creation, negligible genuine post-verification decay in our audited cohort, validated agent-maintenance signals, universal instruction reading, frequent reference following, large behavioral amplification when instructions are present, and no measurable success difference between truthful and confirmed-false referential content under high task difficulty.

Future work must (i) complete aligned ABC coverage on all cases, (ii) pre-register causal load-bearing strata, (iii) replicate across agent platforms, and (iv) separate observational integrity claims from operational-cost claims in reviewer-facing prose.

\paragraph{Data availability.}
All cited metrics are reproducible from the export paths listed in Section~\ref{sec:reproducibility} and \path{exports/paper_synthesis/late_binding_evidence_table.csv}.


\bibliographystyle{ACM-Reference-Format}
\bibliography{references}

\end{document}
